\documentclass{article}
\usepackage{amsmath,amssymb,amsthm}
\usepackage{algorithm}
\usepackage{algpseudocode}
\usepackage{fullpage}
\usepackage{enumitem}
\newtheorem{theorem}{Theorem}
\newtheorem{lemma}{Lemma}
\newtheorem{assumption}{Assumption}
\newtheorem{corollary}{Corollary}
\newcommand{\E}{\mathbb{E}}
\newcommand{\R}{\mathbb{R}}
\newcommand{\norm}[1]{\left\lVert #1\right\rVert}
\newcommand{\ip}[2]{\langle #1 , #2\rangle}
\newcommand{\block}[1]{U_{#1}}
\begin{document}

\section*{Algorithm Name}
\textbf{Randomized Block Coordinate Descent (RBCD)}

\section*{Mathematical Formulation}
Let $f:\R^{d}\rightarrow\R$ be a convex differentiable function.  
Assume that the coordinates are partitioned into $B$ (non‑overlapping) blocks
\[
\{1,\dots ,d\}= \bigcup_{i=1}^{B}\mathcal{B}_{i},\qquad
\mathcal{B}_{i}\cap \mathcal{B}_{j}=\emptyset\;(i\neq j).
\]
For a vector $w\in\R^{d}$ we denote by $w_{\mathcal{B}_{i}}$ the subvector restricted to block $i$ and by
$\nabla_{i}f(w)$ the corresponding block of the gradient.
Let $\block{i}\in\R^{d\times |\mathcal{B}_{i}|}$ be the matrix that extracts block $i$,
i.e. $\block{i}^{\top}w = w_{\mathcal{B}_{i}}$ and $\block{i}\block{i}^{\top}$ zeroes out all coordinates except those in $\mathcal{B}_{i}$.

At iteration $t$ a block $i_{t}\in\{1,\dots ,B\}$ is drawn uniformly at random
($\Pr(i_{t}=i)=1/B$).  With a stepsize $\eta>0$ the update is
\[
\boxed{%
w_{t+1}= w_{t}-\eta\,\block{i_{t}}\nabla_{i_{t}} f(w_{t}) } \tag{1}
\]
Only the selected block is modified; all other coordinates remain unchanged.

\section*{Pseudocode}
\begin{algorithm}[H]
\caption{Randomized Block Coordinate Descent}
\begin{algorithmic}[1]
\State \textbf{Input:} stepsize $\eta>0$, number of blocks $B$, maximum iterations $T$
\State Initialize $w_{0}\in\R^{d}$
\For{$t=0,\dots ,T-1$}
    \State Sample $i_{t}\sim\text{Uniform}\{1,\dots ,B\}$
    \State Compute block gradient $g_{t}= \nabla_{i_{t}} f(w_{t})$
    \State Update only block $i_{t}$:
    \[
        w_{t+1}= w_{t}-\eta\,\block{i_{t}} g_{t}
    \]
\EndFor
\State \textbf{Output:} $w_{T}$
\end{algorithmic}
\end{algorithm}

\section*{Dry Run Example}
Consider the one‑dimensional case (single block) $f(w)=(w-3)^{2}$, $w_{0}=0$, stepsize $\eta=0.1$.

\begin{center}
\begin{tabular}{c|c|c|c}
$t$ & $w_{t}$ & $\nabla f(w_{t})$ & $f(w_{t})$ \\ \hline
0 & $0$ & $2(0-3)=-6$ & $9$ \\
1 & $0-0.1(-6)=0.6$ & $2(0.6-3)=-4.8$ & $(0.6-3)^{2}=5.76$ \\
2 & $0.6-0.1(-4.8)=1.08$ & $2(1.08-3)=-3.84$ & $(1.08-3)^{2}=3.6864$ \\
3 & $1.08-0.1(-3.84)=1.464$ & $2(1.464-3)=-3.072$ & $(1.464-3)^{2}=2.3713$
\end{tabular}
\end{center}
The objective value decreases from $9$ to about $2.37$ after three iterations.

\newpage

\section*{Assumptions}
\begin{assumption}[Convexity]\label{as:convex}
$f$ is convex, i.e. for all $u,v\in\R^{d}$,
\[
f(v)\ge f(u)+\ip{\nabla f(u)}{v-u}.
\]
\end{assumption}

\begin{assumption}[Block‑wise $L_{i}$‑smoothness]\label{as:smooth}
For each block $i\in\{1,\dots ,B\}$ there exists a constant $L_{i}>0$ such that
\[
f\bigl(w+\block{i}d\bigr)\le
f(w)+\ip{\nabla_{i}f(w)}{d}+\frac{L_{i}}{2}\norm{d}^{2},
\qquad\forall\,w\in\R^{d},\; d\in\R^{|\mathcal{B}_{i}|}.
\]
Define $L_{\max}:=\max_{i}L_{i}$.
\end{assumption}

\begin{assumption}[Existence of a minimizer]\label{as:opt}
There exists $w^{\star}\in\arg\min_{w}f(w)$ and the optimal value $f^{\star}=f(w^{\star})$ is finite.
\end{assumption}

\section*{Main Convergence Theorem}
\begin{theorem}[Expected suboptimality of RBCD]\label{thm:rate}
Assume \ref{as:convex}–\ref{as:opt} and let the stepsize be
\[
\eta = \frac{1}{L_{\max}} .
\]
Then for the iterates generated by \eqref{1} we have, for any $T\ge 1$,
\[
\E\bigl[f(w_{T})-f^{\star}\bigr]\;\le\;
\frac{2\,B\,L_{\max}\,\norm{w_{0}-w^{\star}}^{2}}{T}.
\tag{2}
\]
Consequently $ \E[f(w_{T})-f^{\star}]=\mathcal{O}\!\bigl(\tfrac{1}{T}\bigr) $.
\end{theorem}

\section*{Theoretical Properties}
\begin{itemize}[leftmargin=*]
\item \textbf{Stability.} The recursion derived in the proof shows that the distance
$\norm{w_{t}-w^{\star}}^{2}$ is non‑increasing in expectation, guaranteeing stability of the iterates.
\item \textbf{Hyperparameter sensitivity.} The optimal stepsize $\eta=1/L_{\max}$ depends only on the
largest block‑wise Lipschitz constant. Over‑estimation of $L_{\max}$ yields a smaller stepsize and slower convergence, while under‑estimation may break the descent inequality.
\item \textbf{Comparison to full‑gradient GD.} For $B=1$ the bound reduces to the classical $O(1/T)$ rate of gradient descent with stepsize $1/L_{1}$. For $B>1$ the factor $B$ appears, reflecting the fact that only a $1/B$ fraction of coordinates is updated per iteration.
\end{itemize}

\section*{Discussion}
The theorem matches the optimal rate for smooth convex optimization in the stochastic (randomized) block setting.  If, in addition, $f$ is $\mu$‑strongly convex, the same proof yields a linear rate
$\E[f(w_{T})-f^{\star}]\le (1-\frac{\mu}{B L_{\max}})^{T}\bigl(f(w_{0})-f^{\star}\bigr)$.
Thus the analysis presented below is a foundation for both the sublinear and linear regimes.

\newpage

\section*{Preliminaries (Definitions and Lemmas)}
\begin{lemma}[Block smoothness inequality]\label{lem:smooth}
Under Assumption \ref{as:smooth}, for any $w\in\R^{d}$ and any block $i$,
\[
f\!\bigl(w-\eta\,\block{i}\nabla_{i}f(w)\bigr)
\le f(w)-\eta \norm{\nabla_{i}f(w)}^{2}
+\frac{\eta^{2}L_{i}}{2}\norm{\nabla_{i}f(w)}^{2}.
\tag{3}
\]
\end{lemma}
\begin{proof}
Apply the smoothness inequality with $d=-\eta\nabla_{i}f(w)$:
\[
f\bigl(w+\block{i}d\bigr)\le f(w)+\ip{\nabla_{i}f(w)}{d}
+\frac{L_{i}}{2}\norm{d}^{2}.
\]
Since $d=-\eta\nabla_{i}f(w)$, the inner product becomes $-\eta\norm{\nabla_{i}f(w)}^{2}$ and
$\norm{d}^{2}= \eta^{2}\norm{\nabla_{i}f(w)}^{2}$, giving \eqref{3}. ∎
\end{proof}

\begin{lemma}[Expectation over random block]\label{lem:expect}
Let $i_{t}$ be sampled uniformly from $\{1,\dots ,B\}$. For any quantity $Q_{i}$ that depends only on block $i$,
\[
\E_{i_{t}}[Q_{i_{t}}]=\frac{1}{B}\sum_{i=1}^{B} Q_{i}.
\tag{4}
\]
\end{lemma}
\begin{proof}
Uniform sampling implies $\Pr(i_{t}=i)=1/B$ for each $i$, hence the claim follows directly from the definition of expectation. ∎
\end{proof}

\section*{Main Proof (Step‑by‑Step Derivation)}
We prove Theorem \ref{thm:rate}.  Throughout the proof $ \E_{t}[\,\cdot\,] $ denotes
conditional expectation w.r.t. the sigma‑algebra generated by $w_{0},\dots ,w_{t}$,
i.e. the randomness only comes from the choice of $i_{t}$.

\noindent\textbf{Step 1 (Descent inequality for a single block).}  
Apply Lemma \ref{lem:smooth} with stepsize $\eta$ to the randomly selected block $i_{t}$:
\[
f(w_{t+1})\le f(w_{t})-\eta \norm{\nabla_{i_{t}}f(w_{t})}^{2}
+\frac{\eta^{2}L_{i_{t}}}{2}\norm{\nabla_{i_{t}}f(w_{t})}^{2}.
\tag{5}
\]

\noindent\textbf{Step 2 (Take expectation over the random block).}  
Using Lemma \ref{lem:expect} and $L_{\max}\ge L_{i}$ for all $i$,
\begin{align}
\E_{i_{t}}\!\bigl[f(w_{t+1})\mid w_{t}\bigr]
&\le f(w_{t})-\frac{\eta}{B}\sum_{i=1}^{B}\norm{\nabla_{i}f(w_{t})}^{2}
+\frac{\eta^{2}L_{\max}}{2B}\sum_{i=1}^{B}\norm{\nabla_{i}f(w_{t})}^{2} \nonumber\\
&= f(w_{t})-\Bigl(\frac{\eta}{B}-\frac{\eta^{2}L_{\max}}{2B}\Bigr)
\underbrace{\sum_{i=1}^{B}\norm{\nabla_{i}f(w_{t})}^{2}}_{\norm{\nabla f(w_{t})}^{2}} .
\tag{6}
\end{align}
(The equality follows because the block gradients together constitute the full gradient.)

\noindent\textbf{Step 3 (Choose stepsize).}  
Set $\displaystyle \eta=\frac{1}{L_{\max}}$.  Then the coefficient in \eqref{6} simplifies:
\[
\frac{\eta}{B}-\frac{\eta^{2}L_{\max}}{2B}
= \frac{1}{B L_{\max}}-\frac{1}{2B L_{\max}}
= \frac{1}{2B L_{\max}} .
\tag{7}
\]
Hence
\[
\E_{i_{t}}\!\bigl[f(w_{t+1})\mid w_{t}\bigr]
\le f(w_{t})-\frac{1}{2B L_{\max}}\norm{\nabla f(w_{t})}^{2}.
\tag{8}
\]

\noindent\textbf{Step 4 (Relate gradient norm to suboptimality).}  
Convexity (Assumption \ref{as:convex}) gives
\[
f(w_{t})-f^{\star}\le \ip{\nabla f(w_{t})}{w_{t}-w^{\star}}
\le \norm{\nabla f(w_{t})}\,\norm{w_{t}-w^{\star}} .
\tag{9}
\]
Applying Young’s inequality $ab\le \frac{a^{2}}{2\alpha}+\frac{\alpha b^{2}}{2}$ with
$\alpha = 2B L_{\max}$ yields
\[
f(w_{t})-f^{\star}
\le \frac{1}{4B L_{\max}}\norm{\nabla f(w_{t})}^{2}
      +B L_{\max}\norm{w_{t}-w^{\star}}^{2}.
\tag{10}
\]
Rearranging gives a lower bound on the gradient norm:
\[
\norm{\nabla f(w_{t})}^{2}
\ge 4B L_{\max}\bigl(f(w_{t})-f^{\star}\bigr)
    -4B^{2}L_{\max}^{2}\norm{w_{t}-w^{\star}}^{2}.
\tag{11}
\]

\noindent\textbf{Step 5 (Combine \eqref{8} and \eqref{11}).}  
Substituting \eqref{11} into \eqref{8} we obtain
\begin{align}
\E_{i_{t}}\!\bigl[f(w_{t+1})\mid w_{t}\bigr]
&\le f(w_{t})-\frac{1}{2B L_{\max}}
\Bigl(4B L_{\max}\bigl(f(w_{t})-f^{\star}\bigr)
      -4B^{2}L_{\max}^{2}\norm{w_{t}-w^{\star}}^{2}\Bigr) \nonumber\\
&= f(w_{t})-2\bigl(f(w_{t})-f^{\star}\bigr)
   +2B L_{\max}\norm{w_{t}-w^{\star}}^{2} \nonumber\\
&= f^{\star}+2B L_{\max}\norm{w_{t}-w^{\star}}^{2}.
\tag{12}
\end{align}
Taking total expectation (law of total expectation) yields
\[
\E\!\bigl[f(w_{t+1})\bigr]
\le f^{\star}+2B L_{\max}\,\E\!\bigl[\norm{w_{t}-w^{\star}}^{2}\bigr].
\tag{13}
\]

\noindent\textbf{Step 6 (Distance recursion).}  
From the update rule \eqref{1} we have
\[
w_{t+1}-w^{\star}=w_{t}-w^{\star}
               -\eta\,\block{i_{t}}\nabla_{i_{t}}f(w_{t}).
\]
Squaring and taking conditional expectation gives
\begin{align}
\E_{i_{t}}\!\bigl[\norm{w_{t+1}-w^{\star}}^{2}\mid w_{t}\bigr]
&= \norm{w_{t}-w^{\star}}^{2}
   -\frac{2\eta}{B}\ip{\nabla f(w_{t})}{w_{t}-w^{\star}}
   +\frac{\eta^{2}}{B}\norm{\nabla f(w_{t})}^{2}.
\tag{14}
\end{align}
Using convexity \eqref{9} to replace the inner product and the bound
\eqref{11} for $\norm{\nabla f(w_{t})}^{2}$, together with $\eta=1/L_{\max}$,
we obtain
\begin{align}
\E_{i_{t}}\!\bigl[\norm{w_{t+1}-w^{\star}}^{2}\mid w_{t}\bigr]
&\le \norm{w_{t}-w^{\star}}^{2}
   -\frac{2}{B L_{\max}}\bigl(f(w_{t})-f^{\star}\bigr)
   +\frac{1}{B L_{\max}^{2}}
      \bigl(4B L_{\max}(f(w_{t})-f^{\star})\bigr) \nonumber\\
&= \norm{w_{t}-w^{\star}}^{2}
   -\frac{2}{B L_{\max}}\bigl(f(w_{t})-f^{\star}\bigr)
   +\frac{4}{L_{\max}}\bigl(f(w_{t})-f^{\star}\bigr) \nonumber\\
&= \norm{w_{t}-w^{\star}}^{2}
   -\frac{2}{B L_{\max}}\bigl(f(w_{t})-f^{\star}\bigr)
   +\frac{4}{L_{\max}}\bigl(f(w_{t})-f^{\star}\bigr) .
\tag{15}
\end{align}
Since $B\ge 1$, the coefficient of $(f(w_{t})-f^{\star})$ is non‑positive:
\[
-\frac{2}{B L_{\max}}+\frac{4}{L_{\max}}
= \frac{4B-2}{B L_{\max}}\ge 0,
\]
hence we can drop the last term to obtain the simpler bound
\[
\E_{i_{t}}\!\bigl[\norm{w_{t+1}-w^{\star}}^{2}\mid w_{t}\bigr]
\le \norm{w_{t}-w^{\star}}^{2}
   -\frac{2}{B L_{\max}}\bigl(f(w_{t})-f^{\star}\bigr).
\tag{16}
\]

\noindent\textbf{Step 7 (Telescoping sum).}  
Take total expectation of \eqref{16} and sum from $t=0$ to $T-1$:
\[
\E\!\bigl[\norm{w_{T}-w^{\star}}^{2}\bigr]
\le \norm{w_{0}-w^{\star}}^{2}
   -\frac{2}{B L_{\max}}\sum_{t=0}^{T-1}\E\!\bigl[f(w_{t})-f^{\star}\bigr].
\tag{17}
\]
Since the left‑hand side is non‑negative,
\[
\sum_{t=0}^{T-1}\E\!\bigl[f(w_{t})-f^{\star}\bigr]
\le \frac{B L_{\max}}{2}\norm{w_{0}-w^{\star}}^{2}.
\tag{18}
\]

\noindent\textbf{Step 8 (From average to last iterate).}  
The bound \eqref{18} holds for the average of the first $T$ iterates.
Using convexity of $f$ we have
\[
f(w_{T})-f^{\star}
\le \frac{1}{T}\sum_{t=0}^{T-1}\bigl(f(w_{t})-f^{\star}\bigr),
\]
and taking expectation yields
\[
\E\!\bigl[f(w_{T})-f^{\star}\bigr]
\le \frac{1}{T}\sum_{t=0}^{T-1}\E\!\bigl[f(w_{t})-f^{\star}\bigr]
\le \frac{B L_{\max}}{2T}\norm{w_{0}-w^{\star}}^{2}.
\tag{19}
\]

\noindent\textbf{Step 9 (Final constant adjustment).}  
Multiplying numerator and denominator of \eqref{19} by $2$ yields the statement
\eqref{2} of Theorem \ref{thm:rate}:
\[
\E\!\bigl[f(w_{T})-f^{\star}\bigr]
\le \frac{2\,B\,L_{\max}\,\norm{w_{0}-w^{\star}}^{2}}{T}.
\qquad\blacksquare
\]

\section*{Rate Derivation (With Constants)}
The bound \eqref{2} can be written as
\[
\E\!\bigl[f(w_{T})-f^{\star}\bigr]
\le \underbrace{\bigl(2B L_{\max}\norm{w_{0}-w^{\star}}^{2}\bigr)}_{\displaystyle C}
\cdot\frac{1}{T},
\]
where the constant $C$ encapsulates problem‑specific quantities:
\begin{itemize}
\item $B$ – number of blocks (larger $B$ slows convergence proportionally);
\item $L_{\max}$ – worst‑case block smoothness constant;
\item $\norm{w_{0}-w^{\star}}^{2}$ – initialization distance.
\end{itemize}
Thus the algorithm attains the optimal $\mathcal{O}(1/T)$ rate for smooth convex
optimization in expectation.

\section*{Special Cases}
\begin{enumerate}
\item \textbf{Single block ($B=1$).} The method reduces to standard gradient descent with stepsize $1/L_{1}$ and the bound becomes $\displaystyle \E[f(w_{T})-f^{\star}]\le \frac{2L_{1}\norm{w_{0}-w^{\star}}^{2}}{T}$, matching the classic result.
\item \textbf{Strongly convex $f$.} If $f$ satisfies $\mu$‑strong convexity,
\[
f(v)\ge f(u)+\ip{\nabla f(u)}{v-u}+\frac{\mu}{2}\norm{v-u}^{2},
\]
the recursion \eqref{16} yields a linear rate
\[
\E\!\bigl[f(w_{T})-f^{\star}\bigr]
\le \bigl(1-\tfrac{\mu}{B L_{\max}}\bigr)^{T}\!\bigl(f(w_{0})-f^{\star}\bigr).
\]
\end{enumerate}

\end{document}